\chapter{Capacitación}
\label{sec:capacitacion}

\section{Plan de Capacitación}
\label{sec:plan-capacitacion}

El equipo SQA recibió capacitación en las siguientes áreas antes del inicio del sprint:

\begin{longtable}{lp{4cm}p{6cm}}
\caption{Capacitación del equipo SQA}
\label{tab:capacitacion}\\
\toprule
\textbf{Área} & \textbf{Duración} & \textbf{Cobertura}\\
\midrule
\endfirsthead
\toprule
\textbf{Área} & \textbf{Duración} & \textbf{Cobertura}\\
\midrule
\endhead
\bottomrule
\endfoot

Jira Cloud & 2 h & Configuración de proyecto, tablero, issues, workflows\\

SonarCloud & 2 h & Análisis estático, calidad de código, reportes\\

pytest & 3 h & Pruebas unitarias, fixtures, mocks, httpx\\

Cypress & 3 h & Pruebas E2E, comandos personalizados, CI\\

Git & 1 h & Flujo de ramas, resolución de conflictos\\

LaTeX (Overleaf) & 1 h & Edición y compilación de documentos\\

\end{longtable}

\section{Capacitación Requerida}
\label{sec:capacitacion-requerida}

No se identificaron necesidades adicionales de capacitación para el equipo SQA durante este sprint. Para sprints futuros se recomienda capacitación en:

\begin{itemize}
    \item Pruebas de carga con k6 o Locust.
    \item Pruebas de seguridad con OWASP ZAP.
    \item Automatización de accesibilidad con axe-core.
\end{itemize}
